\documentclass[10pt,a4paper]{amsart}
\usepackage[margin=19mm]{geometry}
\usepackage{amsmath,amssymb,mathtools}
\usepackage[scr=rsfs,cal=euler]{mathalfa}
\usepackage{xfrac}
\usepackage[unicode,hidelinks,bookmarks=false]{hyperref}
\newcommand{\ie}{i.e.\ }
\newcommand{\cf}{cf.\ }
\newcommand{\resp}{respectively\ }
\newcommand{\Subsection}{\subsection}
\providecommand{\nobreakdash}{\nobreak\hbox{-}}
\setlength{\emergencystretch}{2em}
\multlinegap0pt
\usepackage{amssymb}
\usepackage{mathtools}
\usepackage[shortlabels]{enumitem}
\newtheorem{lemma}{Lemma}
\newcommand{\F}{\mathcal{F}}
\newcommand{\G}{\mathcal{G}}
\title{Definability from Factorised Symmetry in Ultrapowers\\[2mm] \large (Affine Geometries and Beyond)}
\author{Mike Stannett}
\date{Source: arXiv:2607.29366v1 (CC BY 4.0)}
\begin{document}
\maketitle

\noindent\emph{Source and licence.} Definability from Factorised Symmetry in Ultrapowers\\[2mm] \large (Affine Geometries and Beyond). Version arXiv:2607.29366v1 (CC BY 4.0), distributed by arXiv under the Creative Commons Attribution 4.0 International licence, http://creativecommons.org/licenses/by/4.0/, as recorded in the arXiv licence metadata for this exact version and shown on the abstract page https://arxiv.org/abs/2607.29366v1. Full licence notice: \texttt{licenses/arxiv-2607.29366-CC-BY-4.0.txt}.\par\smallskip

\noindent\textit{Editorial scope.} Counted unit: the article's lemma lem:4.3 (source0.tex lines 450--462). The article's complete original proof of it is delivered with it (source0.tex lines 463--490, one \texttt{proof} environment of 28 lines). No mathematics is written, repaired or re-derived: the statement and the proof are the article's own lines, in the article's own notation, and no retained line is substituted or re-flowed.  No further source-local context is needed: every cross-reference of every retained line resolves inside the two delivered spans.  Nothing was omitted from any retained span, so the package carries no omission record.  No retained line contains a citation, so the wrapper carries no bibliography and no bibliography entry.  No retained line contains a figure environment, an image-inclusion command or drawing code, so no image is delivered and the package declares no asset dependency.  Editorial formatting only, in addition: an AMS \texttt{amsart} preamble that restates the article's own declarations and macros used by the retained lines, namely the environments \texttt{lemma} on separate counters, together with the standard packages those lines need.  The retained environments print in this document as Lemma 1 in order of appearance, whereas the article prints the same items with the numbering of its own full document.  The complete native source of the article as distributed in the arXiv e-print for this exact version is bundled unchanged as \texttt{sources/206477/source0.tex}.
\begin{lemma}[First-order expressibility and transfer of affine
preservation]
\label{lem:4.3}
Let $\G$ be a finite-relational coordinate geometry and
let $R\subseteq (F^d)^n$ be a parameter-free $\F$-definable
relation. The statement
\begin{center}
\emph{every affine automorphism of $\G$ preserves $R$}
\end{center}
is expressible by a first-order sentence in the language of
$\F$ (uniformly in formulas defining the primitives of $\G$
and defining $R$), and hence transfers to ultrapowers.
\end{lemma}

\begin{proof}
	Choose parameter-free $\F$-formulas
$\kappa,\pi_1,\dots,\pi_k$ defining $\mathsf K,P_1,\dots,P_k$, and a
parameter-free formula $\varphi$ defining $R$. For a matrix tuple
$B\in M^{d\times d}$ and a vector $b\in F^d$, write
$T_{B,b}(x)=Bx+b$. The required assertion is expressed by the sentence
\[
\begin{aligned}
\forall B\,\forall b\,\big(&\det(B)\neq 0 \\
&{}\wedge \forall x_1x_2x_3\,
 \big(\kappa(x_1,x_2,x_3)\leftrightarrow \\
&\hspace{43mm}
 \kappa(T_{B,b}x_1,T_{B,b}x_2,T_{B,b}x_3)\Big) \\
&{}\wedge \bigwedge_{i=1}^k \forall \bar x\,
 \big(\pi_i(\bar x)\leftrightarrow\pi_i(T_{B,b}\bar x)\Big) \\
&{}\Longrightarrow \forall \bar x\,
 \big(\varphi(\bar x)\leftrightarrow
 \varphi(T_{B,b}\bar x)\Big)\Bigg).
\end{aligned}
\] Here $T_{B,b}\bar x$ means that $T_{B,b}$ is applied to each
point in the tuple. The finite conjunction over the primitives is
available because $\G$ is finite-relational. The determinant is a fixed
polynomial in the matrix entries and is therefore definable in the field
language. The displayed biconditionals say exactly that the coded affine
bijection is an automorphism of $\G$ and that it preserves $R$.
Transfer to ultrapowers is now an immediate application of Łoś's
theorem.
\end{proof}

\end{document}
